\chapter{Results and Discussion}
\label{ch:results}

This chapter presents the results obtained by the procedures of Chapter~3:
the Foldy--Wouthuysen (FW) mapping from SME coefficients to
Dobrescu--Mocioiu (DM) potentials (Section~\ref{sec:results-fw}), the
compiled experimental constraint database (Section~\ref{sec:results-db}),
the matter--antimatter asymmetry analysis (Section~\ref{sec:results-asym}),
and the resulting coverage gap analysis
(Section~\ref{sec:results-gap}). Each result is discussed as it is
presented, since the physical interpretation of every numerical result in
this thesis depends on the sensitivity-gap caveat established in
Section~\ref{sec:results-fw-bmu} below.

\section{Foldy--Wouthuysen Derivation Results}
\label{sec:results-fw}

\subsection{The \texorpdfstring{$b_\mu$}{b\_mu} Term}
\label{sec:results-fw-bmu}

The CPT-odd term $\mathcal{L}_b = -b_\mu\,\bar\psi\,\gamma_5\gamma^\mu\,\psi$
splits into a spatial part $b_i$, which is even and therefore exact at
order $m^0$, and a temporal part $b_0$, which is odd and enters only at
order $1/m$. Direct computation in the Dirac representation gives
$\gamma^0\gamma_5\gamma^i=-\Sigma^i$, so that, projected onto the upper
block,
\begin{equation}
  H_{NR}(b_i) = -\mathbf b\cdot\boldsymbol\sigma \quad (\text{matter}),
  \label{eq:HNR-bi-matter}
\end{equation}
already in non-relativistic form with no further FW iteration required.
The temporal component is odd, since $\gamma^0\gamma_5\gamma^0=-\gamma_5$
(using $\{\gamma_5,\gamma^0\}=0$ and $(\gamma^0)^2=I$), so it does not
contribute at $m^0$ and instead enters through the FW cross-term between
the free odd operator $\boldsymbol\alpha\cdot\mathbf p$ and the
perturbing odd operator $\mathcal O_{b_0}=-b_0\gamma_5$: expanding
$(\boldsymbol\alpha\cdot\mathbf p+\mathcal O_{b_0})^2$ and keeping only
the piece linear in $b_0$,
\begin{equation}
  H_{NR}(b_0) = \frac{\beta}{2m}\{\boldsymbol\alpha\cdot\mathbf p,\,\mathcal O_{b_0}\}
  \Big|_{\text{upper block}}
  = -\frac{b_0(\boldsymbol\sigma\cdot\mathbf p)}{m} + O(m^{-2}).
\end{equation}
The antiparticle sector follows from the lower block by the hole-theory
reinterpretation. An antiparticle of momentum $\mathbf p$ and spin
$\mathbf s$ is the absence of a negative-energy solution of momentum
$-\mathbf p$ and spin $-\mathbf s$, and its energy is minus the energy of
that solution. A CPT test compares a particle with its CPT-conjugate
state, since CPT takes a particle of momentum $\mathbf p$ and spin
$\mathbf s$ to an antiparticle of the same momentum and spin $-\mathbf s$.
Labelling that antiparticle state by the particle's spin operator, its
Hamiltonian is
\begin{equation}
  H_{NR}^{\bar f}(\mathbf p,\boldsymbol\sigma)
  = -\,H_{\rm lower}(-\mathbf p,\boldsymbol\sigma).
  \label{eq:HNR-anti-rule}
\end{equation}
Since $\Sigma^i=\mathrm{diag}(\sigma^i,\sigma^i)$ repeats the \emph{same}
operator $\sigma^i$ in the upper and lower blocks,
\begin{equation}
  \langle\text{upper}|\,b_i\gamma^0\gamma_5\gamma^i\,|\text{upper}\rangle =
  \langle\text{lower}|\,b_i\gamma^0\gamma_5\gamma^i\,|\text{lower}\rangle =
  -b_i\,\sigma^i,
  \label{eq:HNR-bi-block}
\end{equation}
and because $b_i$ carries no momentum, Eq.~\eqref{eq:HNR-anti-rule} gives
\begin{equation}
  H_{NR}^{\bar f}(b_i) = -\left(-\mathbf b\cdot\boldsymbol\sigma\right)
  = +\mathbf b\cdot\boldsymbol\sigma. \label{eq:HNR-bi-anti}
\end{equation}
The CPT-conjugate antiparticle receives the opposite energy shift to the
particle. This matter--antimatter sign flip is the signature of a CPT-odd
coefficient. It is not a charge-conjugation effect:
the axial-vector bilinear $\bar\psi\gamma_5\gamma^\mu\psi$ is even under $C$
for every $\mu$, so $C$ alone cannot produce the sign difference between
Eqs.~\eqref{eq:HNR-bi-matter} and~\eqref{eq:HNR-bi-anti}. The difference
is a CPT effect, and it appears only when CPT-conjugate states are
compared.

Stated at fixed physical spin, the result looks different. An antiparticle with
the \emph{same} spin as the particle has energy $-\mathbf b\cdot\mathbf s$,
the same as the particle. This is Kostelecký and Lane's antiparticle
rule, which leaves $b_\mu$, $c_{\mu\nu}$ and $g_{\lambda\mu\nu}$ unchanged
and reverses $a_\mu$, $d_{\mu\nu}$, $e_\mu$, $f_\mu$ and $H_{\mu\nu}$
(Kostelecký and Lane, 1999b). The two statements are equivalent because a
spin-linear term changes sign when the spin is reversed. For a CPT test
the CPT-conjugate comparison is the relevant one: with it, and only with
it, exact CPT symmetry gives $g^{\bar f}_\alpha=g^{f}_\alpha$ and
$A_\alpha=0$. We verified both comparisons symbolically for every
coefficient family, using the full Foldy--Wouthuysen reduction and the
kinetic-sector field redefinition. The CPT-conjugate comparison reproduces
each coefficient's CPT parity in every family, and the fixed-spin
comparison reproduces Kostelecký and Lane's rule in every family. Both
comparisons require the momentum reversal in
Eq.~\eqref{eq:HNR-anti-rule}. Omitting it changes the result for every
term that is odd in $\mathbf p$.\footnote{This agrees with Bluhm,
Kosteleck\'y and Russell's (1999) treatment of hydrogen and antihydrogen.
Their substitution rule takes $a_\mu\to-a_\mu$, $d_{\mu\nu}\to-d_{\mu\nu}$
and $H_{\mu\nu}\to-H_{\mu\nu}$ but leaves $b_\mu$ unchanged, which is the
fixed-spin rule above. The $b_\mu$-dependent shift in their 1S--2S
frequency changes sign between the two atoms because ``the hyperfine states
in $\bar{\rm H}$ have opposite positron and antiproton spin assignments
relative to those of the electron and proton in H'': the transitions being
compared connect CPT-conjugate spin states.}
Promoting
Eq.~\eqref{eq:HNR-bi-matter} to a
two-body potential by pairing the same vertex on both fermions gives the
spin--spin structure
\begin{equation}
  V_2(r) = -g_A^{(1)}g_A^{(2)}\,\frac{1}{4\pi}\,
  (\boldsymbol\sigma_1\cdot\boldsymbol\sigma_2)\,\frac{e^{-r/\lambda}}{r},
  \label{eq:V2-final}
\end{equation}
identifying $b_i$ with $V_2$ at leading order, while $b_0$ matches the
velocity-dependent $V_8$ at order $1/m$.

\paragraph{Discussion: the asymmetry parameter does not saturate the way a
naive substitution suggests.} Substituting the exact, signed relation
$g_{\bar f}=-g_f$ implied by Eqs.~\eqref{eq:HNR-bi-matter} and
\eqref{eq:HNR-bi-anti} directly into Eq.~\eqref{eq:asymmetry} gives
\begin{equation}
  A_\alpha = \frac{g-(-g)}{g+(-g)} = \frac{2g}{0} \;\longrightarrow\; \infty,
  \label{eq:Aalpha-divergence}
\end{equation}
a divergence, confirmed by a symbolic limit calculation, not a
value that saturates at $1$. This is not, however, what SPINDEP actually
computes: its $g_f$ and $g_{\bar f}$ inputs are independent,
always-positive experimental \emph{upper bounds}, not a signed measurement
of one underlying coupling. For any two independent positive bounds
$g_1,g_2$, the ratio $(g_1-g_2)/(g_1+g_2)$ approaches $\pm1$ whenever one
bound is far tighter than the other, \emph{regardless of whether the
underlying physics is CPT-odd or CPT-even}: substituting
$g_{\text{tight}}=r\,g_{\text{loose}}$ and taking $r\to0^+$ gives
$A_\alpha\to1$ with no CPT content at all. Consequently, an observed
$|A_\alpha|\approx1$ is \emph{consistent with}, but not \emph{evidence
for}, a non-zero $b_\mu$: the same numerical pattern is exactly what
independent bounds of very different sensitivity would produce even under
exact CPT symmetry. This result governs the interpretation of every
asymmetry value reported later in this chapter.

\subsection{The \texorpdfstring{$H_{\mu\nu}$}{H\_mu\_nu} Term}

$H_{\mu\nu}$ is CPT-even but Lorentz-odd: CPT-conjugate matter and
antimatter states receive identical shifts at tree level, so its
experimental signature is an anisotropy (sidereal variation) rather than a
matter--antimatter asymmetry. Its
magnetic-like spatial components $H_{ij}$ are even and generate, at leading
order $m^0$, a result obtained directly from
$\gamma^0\sigma^{ij}\big|_{\text{upper block}}=\varepsilon^{ijk}\sigma^k$
(verified by explicit $4\times4$ matrix multiplication for all three
independent spatial pairs):
\begin{equation}
  H_{NR}(H_{ij}) = +\boldsymbol{\mathcal H}_B\cdot\boldsymbol\sigma,
  \qquad \mathcal H_B^k=\tfrac12\varepsilon^{ijk}H_{ij},
\end{equation}
matching to the dipole--dipole potential $V_3$. Its electric-like mixed
components $H_{0i}$ are odd, since $\gamma^0\sigma^{0i}=i\gamma^i$ is
block-off-diagonal, and generate, via the same FW cross-term used for
$b_0$ (Section~\ref{sec:results-fw-bmu}) with the perturbing operator
$\mathcal O_{H_E}=i\sum_k H_{0k}\gamma^k$, at order $1/m$,
\begin{equation}
  H_{NR}(H_{0i}) = +\frac1m\,\boldsymbol\sigma\cdot(\mathbf p\times\mathbf H_E),
  \qquad H_E^i\equiv H_{0i},
\end{equation}
matching to the spin--velocity potential $V_7$. Applying
Eq.~\eqref{eq:HNR-anti-rule} gives $H_{NR}^{\bar f}=H_{NR}^f$ for both
components, as their CPT-even character requires. At fixed physical spin
the sign reverses instead, in line with Kostelecký and Lane's rule
$H_{\mu\nu}\to-H_{\mu\nu}$. This coefficient therefore gives
$g^{\bar f}_\alpha=g^f_\alpha$ and $A_\alpha=0$ for the underlying
couplings, in contrast to $b_\mu$.

\subsection{The \texorpdfstring{$d_{\mu\nu}$}{d\_mu\_nu} Term}

Following Kostelecký and Lane (1999a), $d_{\mu\nu}$ is a CPT-even,
traceless, kinetic-sector coefficient entering through a derivative
coupling, $\mathcal L_d = \tfrac12 i\,\bar\psi\,d^{\mu\nu}\gamma_5\gamma_\mu\,
\overleftrightarrow\partial_\nu\,\psi$. Direct computation gives
$\Gamma_0=-\gamma_5$ (odd) and $\Gamma_i=+\Sigma^i$ (even) for
$\Gamma_\mu\equiv\gamma^0\gamma_5\gamma_\mu$, from which
\begin{align}
  H_{NR}(d_{i0}) &= +d_{i0}\,m\,\sigma^i \quad (\text{order } m^1,\text{ mass-enhanced}) \;\to V_2, \\
  H_{NR}(d_{ij}) &= +d_{ij}\,p^j\,\sigma^i \quad (\text{order } m^0 \text{ in velocity}) \;\to V_7, V_8, \\
  H_{NR}(d_{00}) &= +d_{00}\,(\boldsymbol\sigma\cdot\mathbf p) \quad (\text{order } m^0 \text{ in momentum}) \;\to V_8.
\end{align}
All three results match Kostelecký and Lane's own non-relativistic
Hamiltonian (1999a, Eq.~4) exactly, including the signs, once that equation
is read in its own momentum convention. That convention is not stated
explicitly in the text of Eq.~4, which describes $p_j$ as the
three-momentum of the particle, but it is fixed unambiguously by the
free-particle term of the same authors' intermediate relativistic
Hamiltonian (Kostelecký and Lane, 1999b), $m\mathcal{P}_0:=-p_j\gamma^0\gamma^j$. That expression equals
$-\boldsymbol\alpha\cdot\mathbf p$, so reproducing the required free Dirac
Hamiltonian $+\boldsymbol\alpha\cdot\mathbf p+\beta m$ forces their
lower-index $p_j$ to be the covariant component $p_j=-p^j$ rather than the
physical three-momentum. Read that way, the mass-enhanced $d_{i0}$ term and
both momentum-linear terms agree with Eq.~4 term by term; read as physical
momentum, the two momentum-linear terms appear to differ by an overall
sign while the mass term still agrees, which is the signature of a
momentum-index mismatch rather than of a derivation error. The full
comparison, together with an independent implementation of Kostelecký and
Lane's (1999b) field-redefinition procedure, is carried out symbolically in the
supporting notebooks.

\subsection{The \texorpdfstring{$g_{\lambda\mu\nu}$}{g\_lambda\_mu\_nu} Term}
\label{sec:results-fw-g}

$g_{\lambda\mu\nu}$ is CPT-odd and dimensionless, and enters the
kinetic sector through $\tfrac12 g_{\lambda\mu\nu}\sigma^{\lambda\mu}$ in
$\Gamma_\nu$. It is antisymmetric in its first two indices, since
$\sigma^{\lambda\mu}$ is; the third is the derivative index. Only the
$\nu=0$ slice enters $\Gamma_0$, so the same kinetic-sector field
redefinition used for $d_{\mu\nu}$ applies, and the resulting Hamiltonian is
Hermitian.

Writing $\tilde g_j\equiv\tfrac12\varepsilon_{jkl}g_{kl0}$ for the dual of
the spatial--spatial $\nu=0$ block and $(\mathbf g_{00})_k\equiv g_{k00}$,
the reduction separates by index family:
\begin{align}
  H_{NR}(g_{[kl]0}) &= -\,m\,\tilde{\mathbf g}\cdot\boldsymbol\sigma
     \quad (\text{order } m^1,\text{ mass-enhanced}) \;\to V_2, \\
  H_{NR}(g_{[0k]0}) &= +\,\boldsymbol\sigma\cdot(\mathbf p\times\mathbf g_{00})
     \quad (\text{order } m^0) \;\to V_7, \\
  H_{NR}(g_{[kl]j}) &= +\tfrac12\varepsilon_{klm}g_{mlj}\,p^j\,\sigma^k
     \quad (\text{order } m^0 \text{ in velocity}) \;\to V_7, V_8,
\end{align}
with $g_{[0k]j}$ contributing nothing at this order. All three reproduce
Kostelecký and Lane's (1999a, Eq.~4) non-relativistic Hamiltonian term by term, under the
same covariant momentum convention established for $d_{\mu\nu}$ in
Section~\ref{sec:results-fw}. Since that convention was fixed there without
reference to $g_{\lambda\mu\nu}$, its also making this sector agree is an
independent confirmation of the reading.

\paragraph{Matter--antimatter behaviour.} Applying
Eq.~\eqref{eq:HNR-anti-rule} to each family gives an exact sign flip for
all three contributing families, as the CPT-odd character of
$g_{\lambda\mu\nu}$ requires. The momentum reversal in that equation is
essential here. The $g_{[0k]0}$ and $g_{[kl]j}$ terms are linear in
$\mathbf p$, and comparing the lower block at the same momentum instead
would wrongly suggest that these two families do not flip. At fixed
physical spin, all three families give identical particle and antiparticle
shifts, in line with Kostelecký and Lane's rule
$g_{\lambda\mu\nu}\to+g_{\lambda\mu\nu}$. The matter--antimatter signature
is therefore set by the coefficient's CPT parity, for $g_{\lambda\mu\nu}$
as for $b_\mu$.

\subsection{Master SME\texorpdfstring{$\to$}{->}DM Matching Table}
\label{sec:master}

Table~\ref{tab:master} synthesises the four derivations above. Its final
column gives the relation between the particle and antiparticle couplings
for CPT-conjugate states. That relation describes the underlying couplings,
not the published bounds: as Section~\ref{sec:results-fw-bmu} shows, an
$A_\alpha$ computed from one-sided bounds cannot reveal it.

\begin{table}[h]
\centering
\small
\begin{tabular}{@{}l l >{\raggedright\arraybackslash}p{3.3cm} l l >{\raggedright\arraybackslash}p{3cm}@{}}
\toprule
SME coeff. & CPT & NR Hamiltonian & DM potential(s) & $1/m$ order & Matter--antimatter$^\dagger$ \\
\midrule
$b_i$ & Odd & $-\mathbf b\cdot\boldsymbol\sigma$ (matter), $+\mathbf b\cdot\boldsymbol\sigma$ (CPT-conjugate antimatter) & $V_2$ & $m^0$ & $g^{\bar f}=-g^{f}$ \\
$b_0$ & Odd & $-b_0(\boldsymbol\sigma\cdot\mathbf p)/m$ & $V_8$ & $m^{-1}$ & $g^{\bar f}=-g^{f}$ \\
$H_{ij}$ & Even & $+\boldsymbol{\mathcal H}_B\cdot\boldsymbol\sigma$ & $V_3$ & $m^0$ & $g^{\bar f}=g^{f}$ \\
$H_{0i}$ & Even & $+\tfrac1m\boldsymbol\sigma\cdot(\mathbf p\times\mathbf H_E)$ & $V_7$ & $m^{-1}$ & $g^{\bar f}=g^{f}$ \\
$d_{i0}$ & Even & $+d_{i0}\,m\,\sigma^i$ & $V_2$ & $m^1$ (!) & $g^{\bar f}=g^{f}$ \\
$d_{ij}$ & Even & $+d_{ij}\,p^j\,\sigma^i$ & $V_7, V_8$ & $m^0$ (vel.) & $g^{\bar f}=g^{f}$ \\
$d_{00}$ & Even & $+d_{00}\,(\boldsymbol\sigma\cdot\mathbf p)$ & $V_8$ & $m^0$ (mom.) & $g^{\bar f}=g^{f}$ \\
$g_{[kl]0}$ & Odd & $-m\,\tilde{\mathbf g}\cdot\boldsymbol\sigma$ & $V_2$ & $m^1$ & $g^{\bar f}=-g^{f}$ \\
$g_{[0k]0}$ & Odd & $+\boldsymbol\sigma\cdot(\mathbf p\times\mathbf g_{00})$ & $V_7$ & $m^0$ & $g^{\bar f}=-g^{f}$ \\
$g_{[kl]j}$ & Odd & $+\tfrac12\varepsilon_{klm}g_{mlj}p^j\sigma^k$ & $V_7, V_8$ & $m^0$ (vel.) & $g^{\bar f}=-g^{f}$ \\
\bottomrule
\end{tabular}
\caption{Master SME$\to$DM matching table, covering the complete set of SME
coefficients that generate spin-dependent non-relativistic potentials.
Here $\tilde g_j\equiv\tfrac12\varepsilon_{jkl}g_{kl0}$ and
$(\mathbf g_{00})_k\equiv g_{k00}$. $^\dagger$Relation between the particle and antiparticle couplings for
CPT-conjugate states (antiparticle spin reversed), verified symbolically
for every row; it follows the coefficient's CPT parity without exception.
At fixed physical spin every sign in this column reverses (Kostelecký and
Lane, 1999b). Neither relation gives $|A_\alpha|\to1$: a flip makes the
signed $A_\alpha$ diverge (Eq.~\ref{eq:Aalpha-divergence}), and equal
couplings give $A_\alpha=0$. A value near $1$ in the data therefore
reflects a sensitivity gap between independent bounds.}
\label{tab:master}
\end{table}

\section{The Compiled Constraint Database}
\label{sec:results-db}

Applying the compilation methodology of Section~\ref{sec:db-methods}
produced a database of 283 experimental constraint datasets, of which 247
enter the physics analysis. Thirty-six are compiled into the registry but
held out of it, for five separately stated reasons. Twenty-eight are \texttt{Combined} envelope curves, twelve in $g_Ag_V$
and sixteen in $g_pg_s$: aggregates over several experiments, which the
source database labels \texttt{combined} and marks \texttt{review\_only},
so that no single $V_n$ applies to them by construction. One is a $g_pg_p$ astrophysical $e$-$N$ bound that carries no
potential assignment and for which none can be established from the source
database. Two are \texttt{Code-plot-matlab} contributor-template
placeholders, whose coupling label is a source directory named for its file
format rather than its physics content. Five are the 90\% confidence-level
variants of the Cong \textit{et al.} (2025) hydrogen-spectroscopy bounds,
which the source database ships alongside the 95\% CL curves; the two sets
are the same measurement reported at two confidence levels, so counting
both would double-count the experiment and pair each twice against the same
antimatter curve. The 95\% CL set is retained as the more conservative
limit and the only one the source database carries a curation record for.
Table~\ref{tab:coupling-breakdown}
gives both counts by coupling family; every count reported in the remainder
of this chapter is an analysed count unless stated otherwise.

\begin{table}[h]
\centering
\begin{tabular}{@{}lrr@{}}
\toprule
Coupling family & Compiled & Analysed \\
\midrule
$g_Ag_A$ & 69 & 66 \\
$g_pg_s$ & 53 & 37 \\
$g_sg_s$ & 49 & 49 \\
$g_Vg_V$ & 45 & 44 \\
$g_Ag_V$ & 37 & 25 \\
$g_pg_p$ & 17 & 15 \\
$V_1$ (unlabelled coupling) & 11 & 11 \\
\texttt{Code-plot-matlab} (template placeholder) & 2 & 0 \\
\midrule
Total & 283 & 247 \\
\bottomrule
\end{tabular}
\caption{Dataset count by coupling family, as compiled into the registry
and as admitted to the physics analysis.}
\label{tab:coupling-breakdown}
\end{table}

By DM potential, the analysed database is dominated by $V_{4+5}$ (61
datasets), followed by $V_1$ (47), $V_3$ (39), $V_{9+10}$ (31), $V_2$
(20), and $V_{12+13}$ (16), with $V_{1a}$ (14), $V_{11}$ (7), $V_8$ (6),
$V_{2+3}$ (2), $V_{16}$ (2), and $V_{15}$ (2) making up the remainder.
A bare numeric prefix in the source filenames denotes the potential
number itself, so that a \texttt{3}-prefixed file is a bound on $V_3$
alone and only the \texttt{23} prefix denotes the $V_2+V_3$ combination;
every assignment in this table has been validated against the source
database's own curation records, with no disagreement among the 194
datasets those records cover. Two successive
filename-convention passes resolved 26 of the 39 files that initially
carried no recoverable potential token, by giving each an explicit
potential-number prefix the parser could read; the 13 that remained are
among the 36 datasets held out of the analysis above. Of the 247 analysed
datasets, only 22 involve an antimatter sector: 9 in $e$-$\bar p$, 6 in
$e$-$\bar\mu$, 5 in $e$-$e^+$, and one each in $\bar p$-He and the
$dd\mu^+$ molecular ion (both, per Cong \textit{et al.}, 2025, laboratory
bounds on $\bar p$-$N$ and $\mu^+$-$N$ respectively). The six $e$-$\bar\mu$
bounds all come from muonium, which contains an antimuon, so they belong on
the antimatter side for the same reason $dd\mu^+$ does. This imbalance, 22
antimatter datasets against 225 matter-sector ones, is the central
empirical fact the remainder of this chapter works with.

\section{Matter--Antimatter Asymmetry Analysis}
\label{sec:results-asym}

\subsection{The Fifteen Matched Pairs}

SPINDEP's pairing logic (Section~\ref{sec:db-methods}) matched fifteen
matter--antimatter pairs from the analysed database. Table~\ref{tab:pairs}
gives, for each, the mean absolute asymmetry, the autocorrelation-corrected
effective degrees of freedom, the resulting significance, and the 95\%
bootstrap confidence interval.

\begin{table}[p]
\centering
\footnotesize
\begin{tabular}{@{}l l l >{\raggedright\arraybackslash}p{3.4cm} r r c l@{}}
\toprule
Coupling & Potential & Sector & Matter / antimatter source & $\overline{|A_\alpha|}$ & dof$_{\text{eff}}$ & $Z_{\text{eff}}$ & 95\% CI \\
\midrule
$g_Ag_A$    & $V_3$    & $e$-$\bar p$ & Cong \textit{et al.} (2025) / Fadeev \textit{et al.} (2022) & 1.0000 & 8  & 51.1 & [1.0000, 1.0000] \\
$g_Ag_A$    & $V_2$    & $e$-$\bar p$ & Cong \textit{et al.} (2025) / Ficek \textit{et al.} (2018) & 0.9999 & 21  & 34.4 & [0.9999, 0.9999] \\
$g_Ag_A$    & $V_2$    & $e$-$\bar p$ & Karshenboim (2011) / Ficek \textit{et al.} (2018) & 0.9998 & 17  & 38.8 & [0.9998, 0.9998] \\
$g_Ag_A$    & $V_2$    & $e$-$e^+$    & Ficek \textit{et al.} (2017) / Karshenboim (2011) & 0.9892 & 21  & 48.2 & [0.9886, 0.9896] \\
$g_Ag_A$    & $V_3$    & $e$-$e^+$    & Ficek \textit{et al.} (2017) / Fadeev \textit{et al.} (2022) & 0.9539 & 6  & 38.5 & [0.9518, 0.9557] \\
$g_Vg_V$    & $V_3$    & $e$-$e^+$    & Fadeev \textit{et al.} (2022) / Fadeev \textit{et al.} (2022) & 0.9535 & 6  & 37.1 & [0.9516, 0.9552] \\
$g_pg_p$    & $V_3$    & $e$-$e^+$    & Fadeev \textit{et al.} (2022) / Fadeev \textit{et al.} (2022) & 0.9535 & 6  & 37.1 & [0.9516, 0.9552] \\
$g_pg_p$    & $V_3$    & $e$-$\bar p$ & Cong \textit{et al.} (2025) / Ficek \textit{et al.} (2018) & 0.9304 & 6  & 34.3 & [0.9293, 0.9314] \\
$g_Ag_A$    & $V_3$    & $e$-$\bar p$ & Cong \textit{et al.} (2025) / Ficek \textit{et al.} (2018) & 0.9303 & 6  & 34.3 & [0.9292, 0.9313] \\
$g_sg_s$    & $V_1$    & $e$-$e^+$    & Delaunay \textit{et al.} (2017) / Adkins \textit{et al.} (2022) & 0.8727 & 21  & 25.7 & [0.8686, 0.8770] \\
$g_Ag_A$    & $V_3$    & $e$-$\bar p$ & Fadeev \textit{et al.} (2022) / Ficek \textit{et al.} (2018) & 0.8236 & 6  & 37.3 & [0.8168, 0.8299] \\
$g_Ag_A$    & $V_3$    & $e$-$\bar p$ & Fadeev \textit{et al.} (2022) / Fadeev \textit{et al.} (2022) & 0.8044 & 8  & 43.8 & [0.8024, 0.8063] \\
$g_Vg_V$    & $V_3$    & $e$-$\bar p$ & Fadeev \textit{et al.} (2022) / Ficek \textit{et al.} (2018) & 0.7994 & 7  & 31.8 & [0.7971, 0.8015] \\
$g_pg_p$    & $V_3$    & $e$-$\bar p$ & Fadeev \textit{et al.} (2022) / Ficek \textit{et al.} (2018) & 0.7994 & 7  & 31.8 & [0.7971, 0.8015] \\
$g_Ag_A$    & $V_2$    & $e$-$e^+$    & Jiao \textit{et al.} (2020) / Karshenboim (2011) & 0.3336 & 9  & 5.1 & [0.3197, 0.3474] \\
\bottomrule
\end{tabular}
\caption[The fifteen matched matter--antimatter pairs]{All fifteen matched matter--antimatter pairs, with the
autocorrelation-corrected effective degrees of freedom, the weighted-$\chi^2$
significance $Z_{\text{eff}}$ after that correction (one-sided Gaussian
equivalent), and the 95\% bootstrap confidence interval on the mean
asymmetry.}
\label{tab:pairs}
\end{table}

Effective degrees of freedom range from 6 to 21 across the fifteen pairs, a
factor of 14--50 below the 300 grid points, so the autocorrelation
correction of Section~\ref{sec:stat-methods} is far from cosmetic. Every
pair stays significant after it. The weakest is the Jiao/Karshenboim
$e$-$e^+$ pair at $Z_{\text{eff}}=5.1$ ($p\approx2\times10^{-7}$); the
others lie between $Z_{\text{eff}}=26$ and $51$.

Two rows repeat two others. The source compilation files the same Fadeev
\textit{et al.}\ (2022) $V_3$ curves under both $g_Vg_V$ and $g_pg_p$, so
the $g_Vg_V$ and $g_pg_p$ pairs built from them give identical numbers. The
fifteen pairs therefore contain thirteen independent comparisons.

\subsection{Discussion}

Per Section~\ref{sec:results-fw-bmu}, an asymmetry value near 1 in Table
\ref{tab:pairs} is \emph{consistent with}, but not \emph{evidence for}, CPT
violation: a large sensitivity gap between the matter- and
antimatter-sector measurement being compared produces exactly the same
signature regardless of the true underlying physics. This applies
uniformly to every row, including the $g_Ag_A\cdot V_3\cdot e$-$\bar p$
pair at $1.0000$, the highest asymmetry in the database. The
$g_sg_s\cdot V_1\cdot e$-$e^+$ pair at $0.8727$ makes the same point from a
different direction. $V_1$ is spin-independent, so none of the four
coefficients in Table~\ref{tab:master} feeds it (only $a_\mu$, $c_{\mu\nu}$
and $e_\mu$ enter spin-independent terms), yet its asymmetry sits well
inside the range spanned by the $g_Ag_A$ pairs, 0.33 to 1.00. A high
asymmetry is simply what two bounds of very different sensitivity produce,
whatever channel they sit in. The lowest asymmetry in
the table, $0.334$ for the $g_Ag_A\cdot V_2\cdot e$-$e^+$ pair (Jiao
\textit{et al.}, 2020, versus Karshenboim, 2011), is the pair whose
matter- and antimatter-sector bounds are closest in sensitivity among the
fifteen, and is therefore the single most informative data point currently
in the database for constraining $b_\mu$ in the $e$-$e^+$ channel, precisely
because it is \emph{not} saturated near $|A_\alpha|\approx1$.

The four pairs built on the Cong \textit{et al.} (2025) hydrogen bounds
show the effect directly. Each can be set against a pair with the same
coupling, potential and antimatter partner but an older $e$-$p$ bound on
the matter side, and in all four the newer bound gives the higher
asymmetry: 0.9999 against 0.9998 for $g_Ag_A\cdot V_2$ with Ficek
\textit{et al.} (2018), 1.0000 against 0.8044 for $g_Ag_A\cdot V_3$ with
Fadeev \textit{et al.} (2022), and 0.9303 against 0.8236 ($g_Ag_A$) and
0.9304 against 0.7994 ($g_pg_p$) for $V_3$ with Ficek \textit{et al.}
(2018). Improving the matter-sector measurement, with no change whatever
to the antimatter sector or to the underlying physics, drove the asymmetry
\emph{up}. This is the clearest available demonstration
that $A_\alpha$ built from one-sided bounds tracks the sensitivity ratio
rather than CPT status, and it is a direct empirical confirmation of the
analytic argument of Section~\ref{sec:results-fw-bmu}.

\subsection{Figures}

Figures~\ref{fig:atlas-v1}--\ref{fig:atlas-v16} show the constraint atlas
for each of the twelve classified potential labels in the compiled
database. $V_1$ carries two panels: the bounds catalogued as a coupling
$|g|$ (Figure~\ref{fig:atlas-v1}) and those catalogued as the
gravity-normalised Yukawa strength $|\alpha|$
(Figure~\ref{fig:atlas-v1-alpha}), plotted separately because the two
conventions typically differ by more than thirty orders of magnitude and cannot share a y-axis. Several labels, $V_{2+3}$, $V_{4+5}$, $V_{9+10}$, and
$V_{12+13}$, each combine two adjacent catalogue potentials that the
underlying experiments do not distinguish individually. $V_{2+3}$ is the
sparsest of these, with only two datasets: a bound is filed under it only
where the source explicitly reports the $V_2+V_3$ combination, while a
bound on $V_3$ alone is filed under $V_3$. $V_6$, $V_7$, and $V_{14}$ have
no classified dataset at all. Figures~\ref{fig:matter-antimatter}--\ref{fig:matter-antimatter-lowest}
show the matter--antimatter comparison for the highest- and
lowest-asymmetry pairs in Table~\ref{tab:pairs}, and
Figure~\ref{fig:matter-antimatter-v1} shows the $V_1$ pair, included
separately since $V_1$ is, after the filename-convention pass discussed in
Section~\ref{sec:results-db}, now the second most populated potential in
the database.

\clearpage
\begin{figure}[p]
\centering
\includegraphics[width=0.85\textwidth,height=0.9\textheight,keepaspectratio]{constraint_V1.png}
\caption{Constraint atlas for $V_1$.}
\label{fig:atlas-v1}
\end{figure}
\clearpage

\begin{figure}[p]
\centering
\includegraphics[width=0.85\textwidth,height=0.9\textheight,keepaspectratio]{constraint_V1_alpha.png}
\caption{Constraint atlas for $V_1$, gravity-normalised. These bounds are
published as the Yukawa strength $|\alpha|$ relative to gravity rather than
as a coupling $|g|$, and are shown on a separate axis because the two
conventions typically differ by more than thirty orders of magnitude.}
\label{fig:atlas-v1-alpha}
\end{figure}
\clearpage

\begin{figure}[p]
\centering
\includegraphics[width=0.85\textwidth,height=0.9\textheight,keepaspectratio]{constraint_V1a.png}
\caption{Constraint atlas for $V_{1a}$.}
\label{fig:atlas-v1a}
\end{figure}
\clearpage

\begin{figure}[p]
\centering
\includegraphics[width=0.85\textwidth,height=0.9\textheight,keepaspectratio]{constraint_V2.png}
\caption{Constraint atlas for $V_2$.}
\label{fig:atlas-v2}
\end{figure}
\clearpage

\begin{figure}[p]
\centering
\includegraphics[width=0.85\textwidth,height=0.9\textheight,keepaspectratio]{constraint_V3.png}
\caption{Constraint atlas for $V_3$, the best-covered potential in the
antimatter sector (39 datasets, 10 of them antimatter).}
\label{fig:atlas-v3}
\end{figure}
\clearpage

\begin{figure}[p]
\centering
\includegraphics[width=0.85\textwidth,height=0.9\textheight,keepaspectratio]{constraint_V2+3.png}
\caption{Constraint atlas for $V_{2+3}$. Only two datasets carry this
combination label, both matter-sector; the great majority of what earlier
drafts filed under $V_{2+3}$ are $V_3$-only bounds and appear in
Figure~\ref{fig:atlas-v3}.}
\label{fig:atlas-v23}
\end{figure}
\clearpage

\begin{figure}[p]
\centering
\includegraphics[width=0.85\textwidth,height=0.9\textheight,keepaspectratio]{constraint_V4+5.png}
\caption{Constraint atlas for $V_{4+5}$, the most populated potential in the database (61 datasets).}
\label{fig:atlas-v45}
\end{figure}
\clearpage

\begin{figure}[p]
\centering
\includegraphics[width=0.85\textwidth,height=0.9\textheight,keepaspectratio]{constraint_V8.png}
\caption{Constraint atlas for $V_8$.}
\label{fig:atlas-v8}
\end{figure}
\clearpage

\begin{figure}[p]
\centering
\includegraphics[width=0.85\textwidth,height=0.9\textheight,keepaspectratio]{constraint_V9+10.png}
\caption{Constraint atlas for $V_{9+10}$.}
\label{fig:atlas-v910}
\end{figure}
\clearpage

\begin{figure}[p]
\centering
\includegraphics[width=0.85\textwidth,height=0.9\textheight,keepaspectratio]{constraint_V11.png}
\caption{Constraint atlas for $V_{11}$.}
\label{fig:atlas-v11}
\end{figure}
\clearpage

\begin{figure}[p]
\centering
\includegraphics[width=0.85\textwidth,height=0.9\textheight,keepaspectratio]{constraint_V12+13.png}
\caption{Constraint atlas for $V_{12+13}$.}
\label{fig:atlas-v1213}
\end{figure}
\clearpage

\begin{figure}[p]
\centering
\includegraphics[width=0.85\textwidth,height=0.9\textheight,keepaspectratio]{constraint_V15.png}
\caption{Constraint atlas for $V_{15}$.}
\label{fig:atlas-v15}
\end{figure}
\clearpage

\begin{figure}[p]
\centering
\includegraphics[width=0.85\textwidth,height=0.9\textheight,keepaspectratio]{constraint_V16.png}
\caption{Constraint atlas for $V_{16}$.}
\label{fig:atlas-v16}
\end{figure}
\clearpage

\begin{figure}[p]
\centering
\includegraphics[width=0.75\textwidth]{gAgA_V2_ep_Karshenboim2011_vs_Ficek2018.png}
\caption{Matter--antimatter comparison for the $g_Ag_A\cdot V_2\cdot e$-$\bar p$ pair (Karshenboim, 2011, versus Ficek \textit{et al.}, 2018), one of the three highest asymmetries in the compiled database ($\overline{|A_\alpha|}=0.9998$).}
\label{fig:matter-antimatter}
\end{figure}
\clearpage

\begin{figure}[p]
\centering
\includegraphics[width=0.75\textwidth]{gAgA_V2_ee_Jiao2019_vs_Karshenboim2011.png}
\caption{Matter--antimatter comparison for the $g_Ag_A\cdot V_2\cdot e$-$e^+$ pair (Jiao \textit{et al.}, versus Karshenboim, 2011), the lowest asymmetry in the compiled database ($\overline{|A_\alpha|}=0.334$), included alongside Figure~\ref{fig:matter-antimatter} to show that the analysis does not produce a saturated asymmetry for every pair by construction.}
\label{fig:matter-antimatter-lowest}
\end{figure}
\clearpage

\begin{figure}[p]
\centering
\includegraphics[width=0.75\textwidth]{gsgs_V1_ee_Delaunay2017_vs_Adkins2022.png}
\caption{Matter--antimatter comparison for the $g_sg_s\cdot V_1\cdot e$-$e^+$ pair (Delaunay \textit{et al.}, 2017, versus Adkins \textit{et al.}, 2022), the only currently matched $V_1$ pair. $V_1$ is spin-independent, so none of the four spin-dependent coefficients feeds it, yet its asymmetry is comparable to that of the spin-dependent pairs (see the discussion following Table~\ref{tab:pairs}).}
\label{fig:matter-antimatter-v1}
\end{figure}
\clearpage

\section{Gap Analysis}
\label{sec:results-gap}

\subsection{Coverage by Fermion Sector}

Table~\ref{tab:sector-gaps} compares the number of analysed datasets on
each side of the eleven canonical matter--antimatter \emph{sector} pairs
used by SPINDEP. These are pairings of fermion sectors, not the matched
dataset pairs of Section~\ref{sec:results-asym}.

\begin{table}[h]
\centering
\begin{tabular}{@{}l l r r >{\raggedright\arraybackslash}p{3.2cm}@{}}
\toprule
Matter sector & Antimatter sector & Matter & Antimatter & Status \\
\midrule
$e$-$e$     & $e$-$e^+$      & 40 & 5 & Populated (used above) \\
$e$-$p$     & $e$-$\bar p$   & 17 & 9 & Populated (used above) \\
$e$-$\mu$   & $e$-$\bar\mu$  & 0  & 6 & Antimatter only (muonium) \\
$e$-$n$     & $e$-$\bar n$   & 16 & 0 & \textbf{Gap} \\
$\mu$-$\mu$ & $\mu$-$\bar\mu$& 1  & 0 & \textbf{Gap} (thin on both sides) \\
$n$-$p$     & $n$-$\bar p$   & 9  & 0 & \textbf{Gap} \\
$n$-$n$     & $n$-$\bar n$   & 8  & 0 & \textbf{Gap} \\
$p$-$p$     & $p$-$\bar p$   & 0  & 0 & No data either side \\
$e$-$N$     & $e$-$\bar N$   & 50 & 0 & \textbf{Largest gap} \\
$n$-$N$     & $n$-$\bar N$   & 49 & 0 & \textbf{Gap} (second largest) \\
$p$-$N$     & $p$-$\bar N$   & 18 & 0$^\dagger$ & \textbf{Gap} in comparable units \\
\bottomrule
\end{tabular}
\caption{Matter versus antimatter dataset counts per canonical fermion
sector pair. $^\dagger$A $\bar p$-$N$ bound does exist, from the antiprotonic-helium
spectroscopy analysed by Salumbides \textit{et al.} (2014); it is excluded from this count
because it is published as a gravity-normalised strength $|\alpha|$ rather
than a coupling bound $|g|$ and so cannot be compared with the $p$-$N$
datasets without a conversion this framework does not perform. The gap in
\emph{comparable} data is therefore real, but it is a unit-conversion gap
rather than an absence of measurement.}
\label{tab:sector-gaps}
\end{table}

The three nucleus sectors dominate the gap. $e$-$N$ is the single largest
opportunity in the database, with 50 matter-sector datasets against zero in
$e$-$\bar N$, but $n$-$N$ is barely behind it at 49, and $p$-$N$ adds a
further 18; between them these three account for 117 of the 225
matter-sector datasets and carry no antimatter-sector counterpart at all.
The $e$-$n$, $n$-$p$, and $n$-$n$ sectors each have a smaller but still
substantial matter-sector base (8--16 datasets), again with no
antimatter-sector counterpart. The $e$-$\mu$ row is the mirror image: all
six of its bounds come from muonium and so sit on the antimatter side, with
no matter-sector $e$-$\mu$ bound to compare them with.

The eleven sector pairs in Table~\ref{tab:sector-gaps} account for 208
matter-sector and 20 antimatter-sector datasets, 228 in all, of the 247
analysed. (The two antimatter datasets not shown are the $\bar p$-He and
$dd\mu^+$ entries, whose sectors have no counterpart row; this is why the
table's antimatter column totals 20 rather than the 22 quoted in
Section~\ref{sec:results-db}.) The remaining 19 do not belong to a matter--antimatter
sector pair at all: 16 are astrophysical bounds, quoted for a single sector
with no laboratory antimatter counterpart, and the other three are the
$\mu$-$N$, $\bar p$-He and $dd\mu^+$ entries, whose sectors have no
counterpart row in the table.

\subsection{Coverage by Dobrescu--Mocioiu Potential}

Table~\ref{tab:potential-gaps} breaks the same comparison down by DM
potential.

\begin{table}[h]
\centering
\begin{tabular}{@{}l r r >{\raggedright\arraybackslash}p{4.5cm}@{}}
\toprule
Potential & Matter datasets & Antimatter datasets & Status \\
\midrule
$V_{4+5}$  & 58 & 3 & Thin antimatter coverage \\
$V_{9+10}$ & 31 & 0 & \textbf{Gap} \\
$V_3$      & 29 & 10 & Best-covered potential (used throughout Section~\ref{sec:results-asym}) \\
$V_{1a}$   & 14 & 0 & \textbf{Gap} \\
$V_2$      & 17 & 3 & Thin antimatter coverage \\
$V_1$      & 42 & 5 & Thin antimatter coverage \\
$V_{12+13}$& 16 & 0 & \textbf{Gap} \\
$V_{11}$   & 7  & 0 & \textbf{Gap} \\
$V_8$      & 5  & 1 & Thin antimatter coverage \\
$V_{15}$   & 2  & 0 & \textbf{Gap} \\
$V_{16}$   & 2  & 0 & \textbf{Gap} \\
$V_{2+3}$  & 2  & 0 & \textbf{Gap} (combination label) \\
\bottomrule
\end{tabular}
\caption{Matter versus antimatter dataset counts per Dobrescu--Mocioiu potential.}
\label{tab:potential-gaps}
\end{table}

Seven of the twelve classified potentials, $V_{11}$, $V_{15}$, $V_{16}$,
$V_{1a}$, $V_{9+10}$, $V_{12+13}$, and $V_{2+3}$, have \emph{zero}
antimatter-sector datasets. Only $V_3$ has antimatter coverage exceeding a
handful of datasets, which is why fourteen of the fifteen matched pairs in
Section~\ref{sec:results-asym} involve $V_2$ or $V_3$.

Figures~\ref{fig:lambda-coverage-v1}--\ref{fig:matter-antimatter-ratio} show
the dataset inventory and coverage matrix for the same database. They count
\emph{independent bounds} rather than dataset files: 48 of the 247 analysed
datasets are the same measured bound filed under more than one coupling
family, and are collapsed to a single curve before plotting, leaving 199
independent bounds. Row and column totals read off those figures are
therefore systematically smaller than the counts in Tables
\ref{tab:sector-gaps} and~\ref{tab:potential-gaps}, which count files as the
registry records them. The two views agree on every gap; they differ only
in how multiply-filed bounds are counted.

\clearpage
\begin{figure}[p]
\centering
\includegraphics[width=0.85\textwidth,height=0.9\textheight,keepaspectratio]{lambda_coverage_V1.png}
\caption{Dataset inventory for $V_1$.}
\label{fig:lambda-coverage-v1}
\end{figure}
\clearpage

\begin{figure}[p]
\centering
\includegraphics[width=0.85\textwidth,height=0.9\textheight,keepaspectratio]{lambda_coverage_V4p5.png}
\caption{Dataset inventory for $V_{4+5}$, the most populated potential in the database (61 datasets).}
\label{fig:lambda-coverage-v45}
\end{figure}
\clearpage

\begin{figure}[p]
\centering
\includegraphics[width=0.85\textwidth,height=0.9\textheight,keepaspectratio]{lambda_coverage_V9p10.png}
\caption{Dataset inventory for $V_{9+10}$.}
\label{fig:lambda-coverage-v910}
\end{figure}
\clearpage

\begin{figure}[p]
\centering
\includegraphics[width=0.75\textwidth]{pair_coverage_matrix.png}
\caption{Dataset count by potential and sector; antimatter-sector rows are
highlighted. Cells count independent bounds, not dataset files, so totals
fall below those in Tables~\ref{tab:sector-gaps}
and~\ref{tab:potential-gaps}; the gaps identified are the same.}
\label{fig:pair-coverage}
\end{figure}
\clearpage

\begin{figure}[p]
\centering
\includegraphics[width=0.75\textwidth]{matter_antimatter_ratio.png}
\caption{Matter-to-antimatter dataset count ratio by sector and potential.}
\label{fig:matter-antimatter-ratio}
\end{figure}
\clearpage

\subsection{Data-Quality and Framework Limitations}

Two limitations of the current database and framework are distinct from
the physics coverage gaps above and are stated here rather than worked
around silently. First, the 36 datasets held out of the analysis
(Section~\ref{sec:results-db}) are excluded rather than reclassified, and
the reasons differ in kind. Twenty-eight are \texttt{Combined} envelope
curves, twelve in $g_Ag_V$ and sixteen in $g_pg_s$, which the source
database marks \texttt{review\_only} and assigns no potential; these are
not a parsing gap that any filename convention could close, because no
single $V_n$ describes an envelope over several experiments, and they
would have to be decomposed into their constituent bounds before they
could enter a per-potential analysis. The sixteen $g_pg_s$ curves were
identified only when every potential assignment in the database was
validated against the source database's curation records: a local
filename prefix had caused them to be read as $V_{1a}$, which those
records contradict. Two are
contributor-template placeholders, and were previously being read as $V_1$
and $V_2$ datasets on the strength of a trailing digit in their filenames,
which inflated those two potentials by one dataset each. The remaining one
is a $g_pg_p$ astrophysical $e$-$N$ bound: a real measurement, but one whose
potential is unassigned upstream, and assigning it $V_{1a}$ from the
filenames of its siblings in other coupling families would be an inference
rather than evidence. It is held out until the source publication settles
it, and is the one exclusion that a single documentary check could reverse.
The last five are the 90\% CL Cong \textit{et al.} (2025) curves, excluded
as duplicates of the 95\% CL set rather than as defective data; retaining
both would count one experiment twice.
Second, six entries under $V_3$, $V_{9+10}$ and $V_{15}$ report
interaction ranges that are physically implausible, reaching
$\lambda\sim10^{52}$~m against an observable universe of order
$10^{26}$~m, almost certainly the result of a unit token being mis-detected
during compilation. Three of the fifteen matched pairs in
Table~\ref{tab:pairs} draw on an affected dataset: the $g_Ag_A\cdot
V_3\cdot e$-$e^+$ pair on the matter side (Ficek \textit{et al.}, 2017),
and the two $g_Ag_A\cdot V_3\cdot e$-$\bar p$ pairs that use the Ficek
\textit{et al.} (2018) bound on the antimatter side. None of these results
is affected, because the asymmetry is evaluated only on the overlapping
range of the two curves, and in every case the implausible tail lies
outside that overlap: the ranges used are
$[2.5\times10^{-12},1.6\times10^{-5}]$~m,
$[2.0\times10^{-11},1.0\times10^{-8}]$~m and
$[2.0\times10^{-11},1.3\times10^{-8}]$~m. A broader coverage
claim drawn from the full database without this cleanup would nonetheless
not be reliable.

\subsection{Proposed Experimental Strategies}

Three concrete directions follow from the coverage gaps above. First,
future antimatter-sector measurement effort is best prioritised in the
$e$-$N$, $n$-$N$, $p$-$N$, $e$-$n$, $n$-$p$, and $n$-$n$ sectors, which
together account for 150 of the 225 matter-sector datasets in the analysed
database and zero antimatter-sector ones; precision antiproton experiments such as BASE, and
antihydrogen spectroscopy such as ALPHA, are the nearest existing
platforms for producing bounds in these channels. Second, the same
logic applies at the potential level to $V_{1a}$, $V_{9+10}$, and
$V_{12+13}$, each with 14--31 matter-sector datasets and none on the
antimatter side; $V_1$, $V_2$, and $V_8$ have thin but non-zero antimatter
coverage and are better placed as candidates for \emph{improving} existing
precision rather than opening a channel from zero. Third, following the
precedent documented by Cong \textit{et al.} (2025), in which Wu
\textit{et al.} (2023) reanalysed comagnetometer data taken for sidereal
Lorentz- and CPT-violation searches as new bounds on the $V_{12+13}$,
$V_{4+5}$ and $V_{9+10}$ terms, a systematic search of
existing BASE, ALPHA, and comagnetometer publications for data that could
be reinterpreted within this framework is likely to close some of the
identified gap faster, and at lower cost, than commissioning new
measurements from scratch. Across all three directions, a new
antimatter-sector bound is scientifically useful for a CPT test
specifically when it approaches the precision of its matter-sector
counterpart, not merely when it exists: a loose first measurement in an
otherwise empty sector would simply reproduce the same
sensitivity-gap-dominated $|A_\alpha|\approx1$ pattern documented
throughout this chapter, without resolving the CPT question it is meant to
address.
